
\begin{longtable}{ | p{4.0cm} | p{8.4cm} | c | l | }
\hline\rowcolor{lightgray}
名称 & 描述 & 地址 & 访问 \\ \hline
\endfirsthead
\hline
\rowcolor{lightgray}
名称 & 描述 & 地址 & 访问 \\
\hline
\endhead
\endfoot

\multicolumn{4}{|l|}{\textbf{机器模式信息 CSR}} \\\hline
\hyperref[regdesc:MVENDORID]{mvendorid} & 机器模式供应商编号寄存器 & 0xF11 & RO \\\hline
\hyperref[regdesc:MARCHID]{marchid} & 机器模式架构编号寄存器 & 0xF12 & RO \\\hline
\hyperref[regdesc:MIMPID]{mimpid} & 机器模式硬件实现编号寄存器 & 0xF13 & RO \\\hline
\hyperref[regdesc:MHARTID]{mhartid} & 机器模式 hart 编号寄存器 & 0xF14 & RO \\\hline
\multicolumn{4}{|l|}{\textbf{机器模式陷阱设置 CSR}} \\\hline
\hyperref[regdesc:MSTATUS]{mstatus} & 机器模式状态寄存器 & 0x300 & R/W \\\hline
\hyperref[regdesc:MISA]{misa} \footnote{尽管 \hyperref[regdesc:MISA]{misa} 具有读/写属性，但由于它的域是硬连线的，所以写操作无效。在 RISC-V 术语中称为 WARL（写入任意数值读取合法数值）。} & 机器模式 ISA 寄存器 & 0x301 & R/W \\\hline
\hyperref[regdesc:MIDELEG]{mideleg} & 机器模式中断委托寄存器 & 0x303 & R/W \\\hline
\hyperref[regdesc:MIE]{mie} & 机器模式中断使能寄存器 & 0x304 & R/W \\\hline
\hyperref[regdesc:MTVEC]{mtvec} \footnote{\hyperref[regdesc:MTVEC]{mtvec} 仅支持在向量模式下对异常处理进行配置，基地址为 256 字节对齐。} & 机器模式异常向量寄存器 & 0x305 & R/W \\\hline
\multicolumn{4}{|l|}{\textbf{机器模式陷阱处理 CSR}} \\\hline
\hyperref[regdesc:MSCRATCH]{mscratch} & 机器模式暂存寄存器 & 0x340 & R/W \\\hline
\hyperref[regdesc:MEPC]{mepc} & 机器模式异常程序计数器 & 0x341 & R/W \\\hline
\hyperref[regdesc:MCAUSE]{mcause} \footnote{\hyperref[regdesc:MCAUSE]{mcause} 中反映的外部中断 ID 也包括 RISC-V 标准为核心内部中断源预留的 ID。} & 机器模式异常原因寄存器 & 0x342 & R/W \\\hline
\hyperref[regdesc:MTVAL]{mtval} & 机器模式异常值寄存器 & 0x343 & R/W \\\hline
\hyperref[regdesc:MIP]{mip} & 机器模式中断等待寄存器 & 0x344 & R/W \\\hline
\multicolumn{4}{|l|}{\textbf{用户模式陷阱设置 CSR}} \\\hline
\hyperref[regdesc:USTATUS]{ustatus} & 用户模式状态寄存器 & 0x000 & R/W \\\hline
\hyperref[regdesc:UIE]{uie} & 用户模式中断使能寄存器 & 0x004 & R/W \\\hline
\hyperref[regdesc:UTVEC]{utvec} & 用户模式异常向量寄存器 & 0x005 & R/W \\\hline
\multicolumn{4}{|l|}{\textbf{用户模式陷阱处理 CSR}} \\\hline
\hyperref[regdesc:USCRATCH]{uscratch} & 用户模式暂存寄存器 & 0x040 & R/W \\\hline
\hyperref[regdesc:UEPC]{uepc} & 用户模式异常程序计数器 & 0x041 & R/W \\\hline
\hyperref[regdesc:UCAUSE]{ucause} & 用户模式异常原因寄存器 & 0x042 & R/W \\\hline
\hyperref[regdesc:UIP]{uip} & 用户模式中断等待寄存器 & 0x044 & R/W \\\hline
\multicolumn{4}{|l|}{\textbf{物理存储保护 (PMP) CSR}} \\\hline\hyperref[sec:riscvcpu-pmp-reg-desc]{pmpcfg0} & 物理存储器保护配置寄存器 & 0x3A0 & R/W \\\hline
\hyperref[sec:riscvcpu-pmp-reg-desc]{pmpcfg1} & 物理存储器保护配置寄存器 & 0x3A1 & R/W \\\hline
\hyperref[sec:riscvcpu-pmp-reg-desc]{pmpcfg2} & 物理存储器保护配置寄存器 & 0x3A2 & R/W \\\hline
\hyperref[sec:riscvcpu-pmp-reg-desc]{pmpcfg3} & 物理存储器保护配置寄存器 & 0x3A3 & R/W \\\hline
\hyperref[sec:riscvcpu-pmp-reg-desc]{pmpaddr0} & 物理存储器保护地址寄存器 & 0x3B0 & R/W \\\hline
\hyperref[sec:riscvcpu-pmp-reg-desc]{pmpaddr1} & 物理存储器保护地址寄存器 & 0x3B1 & R/W \\\hline
\multicolumn{4}{|c|}{\textbf{……}} \\\hline
\hyperref[sec:riscvcpu-pmp-reg-desc]{pmpaddr15} & 物理存储器保护地址寄存器 & 0x3BF & R/W \\\hline
\multicolumn{4}{|l|}{\textbf{触发器模块 CSR（与调试模块共用）}} \\\hline
\hyperref[regdesc:TSELECT]{tselect} & 触发器选择寄存器 & 0x7A0 & R/W \\\hline
\hyperref[regdesc:TDATA1]{tdata1} & 触发器抽象数据寄存器 1 & 0x7A1 & R/W \\\hline
\hyperref[regdesc:TDATA2]{tdata2} & 触发器抽象数据寄存器 2 & 0x7A2 & R/W \\\hline
\hyperref[regdesc:TCONTROL]{tcontrol} & 全局触发器控制寄存器 & 0x7A5 & R/W \\\hline
\multicolumn{4}{|l|}{\textbf{调试模式 CSR}} \\\hline
\hyperref[regdesc:DCSR]{dcsr} & 调试模式控制与状态寄存器  & 0x7B0 & R/W \\\hline
\hyperref[regdesc:DPC]{dpc} & 调试 PC 寄存器 & 0x7B1 & R/W \\\hline
\hyperref[regdesc:DSCRATCH0]{dscratch0} & 调试暂存寄存器 0 & 0x7B2 & R/W \\\hline
\hyperref[regdesc:DSCRATCH1]{dscratch1} & 调试模式暂存寄存器 1 & 0x7B3 & R/W \\\hline
\multicolumn{4}{|l|}{\textbf{性能计数 CSR （自定义）\footnote{这些自定义机器模式 CSR 已经在 RISC-V 标准为用户保留的地址空间中实现。}}} \\\hline
\hyperref[regdesc:MPCER]{mpcer} & 机器模式性能计数器事件寄存器 & 0x7E0 & R/W \\\hline
\hyperref[regdesc:MPCMR]{mpcmr} & 机器模式性能计数器模式寄存器 & 0x7E1 & R/W \\\hline
\hyperref[regdesc:MPCCR]{mpccr} & 机器模式性能计数器计数寄存器 & 0x7E2 & R/W \\\hline
\multicolumn{4}{|l|}{\textbf{GPIO 访问 CSR（自定义）}} \\\hline
\hyperref[regdesc:GPIO_OEN]{cpu\_gpio\_oen} & GPIO 输出使能寄存器 & 0x803 & R/W \\\hline
\hyperref[regdesc:GPIO_IN]{cpu\_gpio\_in} & GPIO 读输入值寄存器 & 0x804 & RO \\\hline
\hyperref[regdesc:GPIO_OUT]{cpu\_gpio\_out} & GPIO 写输出值寄存器 & 0x805 & R/W \\\hline
\multicolumn{4}{|l|}{\textbf{物理存储器属性检查器 (PMAC) CSR}} \\\hline
\hyperref[sec:riscvcpu-pma-reg-desc]{pma\_cfg0} & 物理存储器属性配置寄存器 & 0xBC0 & R/W \\\hline
\hyperref[sec:riscvcpu-pma-reg-desc]{pma\_cfg1} & 物理存储器属性配置寄存器 & 0xBC1 & R/W \\\hline
\hyperref[sec:riscvcpu-pma-reg-desc]{pma\_cfg2} & 物理存储器属性配置寄存器 & 0xBC2 & R/W \\\hline
\hyperref[sec:riscvcpu-pma-reg-desc]{pma\_cfg3} & 物理存储器属性配置寄存器 & 0xBC3 & R/W \\\hline
\multicolumn{4}{|c|}{\textbf{……}} \\\hline
\hyperref[sec:riscvcpu-pma-reg-desc]{pma\_cfg15} & 物理存储器属性配置寄存器 & 0xBCF & R/W \\\hline
\hyperref[sec:riscvcpu-pma-reg-desc]{pma\_addr0} & 物理存储器属性地址寄存器 & 0xBD0 & R/W \\\hline
\hyperref[sec:riscvcpu-pma-reg-desc]{pma\_addr1} & 物理存储器属性地址寄存器 & 0xBD1 & R/W \\\hline
\multicolumn{4}{|c|}{\textbf{……}} \\\hline
\hyperref[sec:riscvcpu-pma-reg-desc]{pma\_addr15} & 物理存储器属性地址寄存器 & 0xBDF & R/W \\\hline

\end{longtable}
